\scanpage{52}

Algèbres étales et le différente (4)\hfill Tribu n° 65

§ 2. \uline{Différente.}

\uline{n° 1.}\quad \uline{Idéaux de \unsure{Filting|Fitting} d'un module} $M$.

Voir thèse de D. Lazard. Lien avec le contenu. Le premier \unsure{Filting|Fitting} est compris entre $\operatorname{Ann}(M)$ et une de ses puissances.

\uline{n° 2.}\quad \uline{Différente de Kähler.}

On prend pour $B$ une algèbre de p. f. sur $A$ (plus généralement il suffit peut être de supposer que $\Omega_A(B)$ est un $B$-module de p. f.). Par définition la différente de Kähler $\uline{D}_{B/A}$ est le premier idéal de \unsure{Filting|Fitting} de $\Omega_A(B)$ :
\margin{\unsure{!}}

si
\diagram{\begin{tikzcd}[ampersand replacement=\&,column sep=small]
A^p\arrow[r,"u"]\&A^q\arrow[r]\&\Omega_A(B)\arrow[r]\&0
\end{tikzcd}}
, c'est $\operatorname{Im}(\Lambda^q u)$.

La variété des zéros de $\uline{D}_{B/A}$ est l'ensemble des points de ramification. Comportement de $\uline{D}_{B/A}$ par changement d'anneau (c'est trivial).

\uline{n° 3.}\quad \uline{Transitivité.}

Dire ce qu'on sait. Si on a $A\to B\to C$ avec $C$ intersection complète locale sur $B$, la formule de transitivité marche.

\uline{n° 4.}\quad \uline{Lien avec la trace.}

Définir la différente de Dedekind par la trace, et montrer qu'elle coïncide avec celle de Kähler quand $A$ et $B$ sont des anneaux de valuation discrète (il s'agit du gros th. 3 du n° 7, § 2, avec son dévissage, -- éventuellement agrémenté des lemmes du n° 6 s'ils n'ont pas été faits dans un chapitre antérieur pour illustrer la notion d'intersection complète. En tous cas on aura dit que si $A$ et $B$ sont réguliers, $A\to B$ est une intersection complète relative. La prop. 11 sera éclairée par le lemme suivant :

``Si $k\to R$ est finie et si $R/\operatorname{rad}(R)$ est étale, alors $R$ monogène ssi les idéaux maximaux de $R$ sont homogènes''.

La prop. 12 est absorbée par le résultat ci-dessus sur les anneaux réguliers.

\uline{n° 5.}\quad \uline{Compléments.}

La ramification modérée du n° 5 (noter que la formule de réduction des traces utilisée en haut de la p. 25 est faite en Alg. VIII, § 11).

Donner la description Kummerienne des extensions totalement et modérément ramifiées (extraire la racine n-ième d'une uniformisante bien choisie ; ceci dans le cas complet); exemple des surfaces de Riemann (ici le groupe de Galois a un générateur canonique $\exp\left(\dfrac{2i\pi}{n}\right)$).

\mspage{51}
